% from ega4-15-fr.tex
\oldpage[IV-3]{229}
pour tous les points d'un voisinage de $x$ résulte de (12.1.1, (vii)).
L'assertion (ii) résulte donc de ce qui précède et de (3.3.4).
\end{proof}

\begin{proposition}[15.3.2]
\label{IV.15.3.2-fr}
Les notations étant celles de (15.3.1), supposons que les conditions (i), (ii),
(iii) de (15.2.2) soient vérifiées, et de plus que $f^{-1}(y)$ soit
équidimensionnel au point $x$. Alors $f$ est équidimensionnel au point $x$.
\end{proposition}

\begin{proof}
On a, pour tout $y'\in Y$, $\operatorname{Supp}(\mathcal F_{y'})=f^{-1}(y')$.
Par hypothèse $\mathcal F_y$ est réduit (donc vérifie $(S_1)$) et est
équidimensionnel au point $x$. En vertu de (12.1.1, (ii)), on peut donc supposer
que $\mathcal F_{f(x')}$ est équidimensionnel et de dimension constante
$e=\dim_x(f^{-1}(y))$ au point $x'$ pour tous les $x'\in U$, ce qui signifie que
$f^{-1}(f(x'))$ est équidimensionnel et de dimension $e$ au point $x'$. Comme
par ailleurs $\mathcal F|U$ est $f$-plat, il résulte de (2.3.4) et de
(13.3.1, $a$)) que $f$ est équidimensionnel au point $x$.
\end{proof}

\begin{corollary}[15.3.3]
\label{IV.15.3.3-fr}
Les notations étant celles de (15.3.1), supposons vérifiées les conditions
suivantes~:
\begin{enumerate}[label=(\roman*)]
  \item $\operatorname{Supp}(\mathcal F)=X$.
  \item $f$ est universellement ouvert aux points génériques des composantes
  irréductibles de $f^{-1}(y)$ contenant $x$.
  \item $\mathcal F_y$ est géométriquement ponctuellement intègre sur
  $\mathbf{k}(y)$ au point $x$ (4.6.22).
  \item $Y$ est géométriquement unibranche au point $y$.
\end{enumerate}
Alors $x$ n'appartient qu'à une seule composante irréductible de $X$.

Si de plus $Y$ est intègre au point $y$ et $\mathcal F=\mathcal O_X$, alors $X$
est intègre au point $x$ et $f$ est plat au point $x$.
\end{corollary}

\begin{proof}
Notons d'abord que (i) et (iii) entraînent que $x$ n'appartient qu'à une seule
composante irréductible $Z$ de $f^{-1}(y)$. Il résulte donc de (14.4.7) et de
(15.3.2) que pour toute composante irréductible $X_i$ de $X$ contenant $x$
(donc $Z$), la restriction $f_i$ de $f$ à $X_i$ est un morphisme
équidimensionnel au point $x$, et universellement ouvert au point générique de
$Z$. La première assertion de l'énoncé résulte alors de (15.1.3) et la seconde
de (15.3.1).
\end{proof}

\begin{remarks}[15.3.4]
\label{IV.15.3.4-fr}
\begin{enumerate}[label=(\roman*)]
  \item L'exemple (14.4.10, (ii)) montre que, dans (15.3.3), on ne peut
  supprimer l'hypothèse que $Y$ est géométriquement unibranche au point $y$.
  \item La remarque (15.2.4, (vi)) montre que la conclusion de (15.3.3) est
  encore vraie sous les hypothèses suivantes~: $\mathcal O_y$ est régulier,
  $\operatorname{Supp}(\mathcal F)=X$ et $\mathcal F_y$ est intègre au point
  $x$. Nous ignorons si dans cet énoncé, on peut remplacer l'hypothèse que
  $\mathcal O_y$ est régulier par celle qu'il est géométriquement unibranche, ou
  même intègre et intégralement clos.
\end{enumerate}
\end{remarks}

\subsection{Compléments sur les morphismes de Cohen-Macaulay}
\label{subsection:IV.15.4-fr}

\begin{proposition}[15.4.1]
\label{IV.15.4.1-fr}
Soient $Y$ un préschéma localement noethérien, $f:X\to Y$ un morphisme localement
de type fini, $\mathcal F$ un $\mathcal O_X$-Module cohérent tel que
$\operatorname{Supp}(\mathcal F)=X$, $x$ un point de $X$, $y=f(x)$. On pose,
pour tout $z\in X$, $X_{f(z)}=f^{-1}(f(z))$,
$\mathcal F_{f(z)}=\mathcal F\otimes_{\mathcal O_{f(z)}}\mathbf{k}(f(z))$.
Supposons que $\mathcal F$ soit $f$-plat au point $x$ et que $\mathcal F_y$ soit
un $\mathcal O_{X_y}$-Module de Cohen-Macaulay au point $x$. Alors $f$ est
